
\subsubsection{LLM Repair Loops with Execution Feedback}

The execution-feedback repair paradigm was established by Self-Debug \citep{Chen2023Teaching}, which demonstrates that providing LLMs with their own program output --- captured stdout, stderr, and exception traces --- enables iterative repair that substantially improves pass@k over single-attempt generation. Self-Debug operates on execution-only feedback, treating the binary pass/fail signal augmented with error output as sufficient for repair. This becomes the Condition A baseline in this work.

Reflexion \citep{Shinn2023Reflexion} extends the repair paradigm by having the LLM generate verbal reflections on its failures before attempting repair. Verbal reflection is unstructured: the LLM authors its own diagnostic narrative. The present work tests a complementary hypothesis --- that external structured diagnostics (mypy output) provide signal beyond what the LLM can self-generate or extract from execution output.

Olausson et al. [2024] find that self-repair gains are often modest when API cost is accounted for and that repair effectiveness varies substantially within and between datasets. ''How Many Tries Does It Take?'' (arxiv 2604.10508) reports that execution-only repair improves HumanEval pass@1 by 4.9 to 17.1 percentage points and MBPP by 16.0 to 30.0 percentage points across seven models with up to five rounds, establishing the execution-only baseline landscape.

\textbf{Gap addressed:} Prior repair loop work treats feedback channel selection as an engineering choice, not an experimental variable. Self-Debug, Reflexion, and their successors do not ablate execution versus execution+static-analysis. This work provides the first controlled comparison.

\subsubsection{Static Analysis in Code Generation and Repair}

Static analysis tools have been integrated into LLM code generation systems without controlled ablation of their marginal contribution. SWE-agent \citep{Yang2024SWEagent} uses a rich tool-use environment including bash execution, file editing, and code analysis, embedded in an agent control loop. The agent can invoke linters, but the bundled design makes it difficult to isolate the contribution of any single feedback channel.

LLMloop (arxiv 2603.23613) integrates compilation errors and static analysis (PMD, Java) in an iterative loop, finding that static analysis adds signal beyond execution-only in that setting. The architecture is Java-specific and uses a different static analysis tool; the present work tests the analogous hypothesis for Python with mypy on standard Python code generation benchmarks.

Automated program repair (APR) surveys \citep{Xia2023Automated} document that execution-based feedback is the dominant repair signal in the literature. Type-system feedback appears in some systems but controlled comparison against execution-only baselines on Python benchmarks is absent.

\textbf{Gap addressed:} Static analysis is routinely bundled with execution feedback in production systems, but its isolated marginal contribution has not been measured on standard Python code generation benchmarks.

\subsubsection{Test-Based Repair as Alternative Structured Feedback}

An alternative approach to providing structured feedback uses LLM-generated test cases. ChatUniTest and CodaMosa generate additional tests to guide repair beyond execution failure on fixed test suites. LLM-generated tests represent a different structured signal --- input-output examples --- rather than diagnostic error messages. These approaches and static analysis feedback are not mutually exclusive; the present work restricts attention to static analysis to isolate its individual contribution.

\subsubsection{Evaluation Benchmarks}

EvalPlus \citep{Liu2023EvalPlus} introduces HumanEval+ and MBPP+ --- enhanced versions of HumanEval \citep{Chen2021Evaluating} and MBPP \citep{Austin2021Program} with substantially more test cases per problem. HumanEval+ (164 problems) and MBPP+ (378 problems) are the primary evaluation benchmarks in this work. A structural observation of this work is that HumanEval+ and MBPP+ have dramatically different mypy error profiles in GPT-4o-mini failing solutions (70\% vs 0\%). Prior work does not report mypy error rates by benchmark, leaving this divergence unobserved.

\subsubsection{Formal Verification for Code Correctness}

Recent work explores whether LLMs can generate formal specifications that constraint solvers then check. The neural-symbolic program synthesis literature (e.g., Rosette \citep{Torlak2014Growing}) uses SMT solvers for specification-guided synthesis; the present work uses LLM-generated specifications, trading completeness for automation. The h-z1 sub-hypothesis tests whether LLM-generated Z3 specs can provide useful counterexample feedback for repair, finding that while spec validity is high (91.1\%), the counterexample activation rate is low (16.1\%) on the arithmetic HumanEval+ subset.

